\setlength{\textfloatsep}{12pt plus 2pt minus 2pt}
\setlength{\intextsep}{10pt plus 2pt minus 2pt}
\section{Introduction}\label{sec:intro}\label{sec:intro-background}

Multiplying two $n$-bit integers is a basic problem in the theory of computation. In the standard deterministic multitape Turing machine model, a fixed finite-state control operates a fixed number of one-dimensional tapes, and the running time counts every head step. In 1971, Sch\"onhage and Strassen gave an $O(n\log n\log\log n)$ multiplication algorithm and conjectured that $n\log n$ steps are required up to constant factors. Harvey and van der Hoeven reached the $O(n\log n)$ bound in 2021. A 2026 preprint then gave an algorithm below $n\log n$, contradicting the conjecture; its construction is intricate, and simplifying it is the subject of this paper.

After Sch\"onhage and Strassen~\cite{SS71}, F\"urer~\cite{Furer07,Furer09} replaced the $\log\log n$ factor by $2^{O(\log^* n)}$, with a modular variant by De, Kurur, Saha, and Saptharishi~\cite{DKSS13}. Harvey, van der Hoeven, and Lecerf~\cite{HvdHL16} sharpened the base of that exponent, Harvey and van der Hoeven~\cite{HvdH21} achieved $O(n\log n)$, and a recent preprint~\cite{OAI26} obtained $O\bigl(n(\log n)^{1-\kappa}\bigr)$ with $\kappa=2^{-182}$. That power saving comes from two small fixed linear networks, yet the algorithm runs them inside the full architecture of Harvey and van der Hoeven. Specifically, it splits the input across several prime-size grids, interpolates onto power-of-two grids, and computes the central transform through polynomial arithmetic. All of that surrounding machinery exists only to connect the two networks to integer multiplication. It accounts for most of the algorithmic complexity and most of the loss in $\kappa$; Section~\ref{sec:overview-barrier} gives the precise accounting.

In this paper, we keep the two finite linear networks of~\cite{OAI26} and replace everything around them by a construction on a single power-of-two grid. The construction is simple enough that every step can be implemented in exact arithmetic and tested bit for bit, and the whole proof verified formally in Lean~4. The $O(n\log n)$ algorithm of Harvey and van der Hoeven takes dimension 1729 in its proof and, to our knowledge, has never been executed.

First, our algorithm slices both $n$-bit inputs into short binary digits and packs them into one array of complex numbers. Second, it arranges that array on one multidimensional power-of-two grid and smooths each axis after the first, called a \emph{spread axis}, using one data-independent localized kernel per spread axis. Third, it runs one discrete Fourier transform directly on that grid using the two networks, and undoes the smoothing to working precision by one two-term cyclic pass per spread axis. Fourth, it squares the transformed values pointwise, runs the same three-step transform pass once more, rounds to integers, and propagates carries to output the exact $2n$-bit product. Algorithm~\ref{alg:main} is stated in full in Section~\ref{sec:overview-five-steps}. As a by-product, the exponent improves to $\kappa=2^{-46}$.\looseness=-1

\subsection{Main results}\label{sec:intro-results}

Our primary algorithmic theorem states the $O\bigl(n(\log n)^{1-2^{-46}}\bigr)$ running-time bound and records the layers that the outer reduction removes.

\begin{maintheorem}\label{thm:main}
Algorithm~1 (Section~\ref{sec:overview-five-steps}), implemented on a deterministic multitape Turing machine with a fixed finite alphabet and a fixed finite number of tapes, multiplies two $n$-bit integers exactly in worst-case time
\[
  O\bigl(n(\log n)^{1-2^{-46}}\bigr).
\]
\emph{Structure.} The outer algorithm uses only power-of-two DFTs on one grid, one fixed convolution kernel per spread axis, and one two-term cyclic recurrence per spread axis; it uses no Chinese remainder map, prime transform length, coprime-length resampling, Fourier-grid expansion, Bluestein chirp, or polynomial-ring convolution. Ordinary products of $O((\log n)^2)$ bits (and inputs below a fixed threshold $n_*$) are delegated to any independent multiplier with $\mathsf M(N)=N(\log N)^{O(1)}$ (for instance that of Sch\"onhage and Strassen~\cite{SS71}), with no recursive call to Algorithm~1 itself.
\end{maintheorem}

Our second main theorem gives a complete formal proof of this complexity bound in the Lean~4 theorem prover without importing any external network bound.

\begin{maintheorem}[Formally verified bound]\label{thm:lean}
The $O\bigl(n(\log n)^{1-2^{-46}}\bigr)$ multiplication bound of Theorem~\ref{thm:main} is formally proved in Lean~4~\cite{dMU21,Mathlib20} in Mathlib's finite-alphabet multi-stack Turing machine model over~$\N$ (statement \texttt{SubNLogN.Challenge.IntMulSubNLogN} in~\cite{Lean26}, available in our \href{https://github.com/lockqueue/intmul}{github repository}; Section~\ref{sec:lean}), depending only on Lean's three standard logical axioms (\texttt{propext}, \texttt{Classical.choice}, \texttt{Quot.sound}) and assuming no result of~\cite{OAI26}.
\end{maintheorem}

\subsection{Executability, formal verification, and use of AI tools}\label{sec:intro-verification}

Because asymptotic multiplication proofs combine rounding budgets with sequential tape invariants, we pair the written proof with an exact-arithmetic reference package and a complete formalization.

\paragraph{Executability.}
Every step of Algorithm~\ref{alg:main} is implemented in exact arithmetic in Python~3.13, with no floating point on any data path of the multiplier or certifier~\cite{Code26}. Across $2{,}067$ tests, $1{,}832$ pass and one is deferred in the fast suite while $234$ pass in the slow suite, checking every returned product bit for bit. All $71$ transform runs with $d\in\{2,\dots,5\}$ and $R\le 1024$ stay within our proved bounds. With its inner-multiplier hierarchy and no stub, the pipeline ran end to end on $70$- and $96$-bit inputs under a test schedule forcing the sub-$n\log n$ route with small substitute networks in place of the $h=100$ networks. The unforced schedule selects the two-dimensional multiplier at every feasible size.\looseness=-1

The $h=100$ network certificates, including the $140$-digit guard constant, are recomputed in exact arithmetic; the $h=100$ networks themselves are certified, not executed. The package is a reference for bit-exact checking, not a fast multiplier, and no running time is measured or claimed. Table~\ref{tab:executability} in Appendix~\ref{app:impl} summarizes the package in our \href{https://github.com/lockqueue/intmul}{github repository}; our companion paper~\cite{MS} develops and formally verifies the two-dimensional instance. As the sub-$n\log n$ regime begins only at astronomical sizes, no practical speedup is claimed.\looseness=-1

\paragraph{Formal verification.}
To remove doubt about hidden hypotheses or rounding-budget gaps, we formalized the $\kappa=2^{-46}$ complexity bound in Lean~4, as detailed in Section~\ref{sec:lean}. To our knowledge, this is the first formally verified proof of an $o(n\log n)$ integer-multiplication bound, as well as the first independent formal verification of the $h=100$ finite-network certificates of~\cite{OAI26}. Building on Mathlib, the development spans roughly 299{,}000 lines across 1{,}062 files for the sub-$n\log n$ proof. It also includes 949 files and about 185{,}000 lines formalizing the two-dimensional $O(n\log n)$ multiplier of our companion paper~\cite{MS}, which instantiates the delegated inner multiplier $\mathsf M$. That inner formalization has headline theorem \texttt{IntMulTM.\allowbreak intMul\_\allowbreak TM2\_\allowbreak nlogn} and is kernel-checked within the same archive, whereas that companion paper distributes a separate Lean development with headline \texttt{IntMul.\allowbreak tmIntMul\_\allowbreak evalsToInTime\_\allowbreak n\_\allowbreak log\_\allowbreak n}.\looseness=-1

Our formal development~\cite{Lean26} exports one closed theorem. We replayed it through the Lean kernel with \texttt{leanchecker} and checked it again with two independent external checkers, \texttt{comparator} and \texttt{nanoda}. What remains to be trusted is small: Lean's kernel, the Mathlib definitions named in the statement, the three standard axioms, and the 80-line statement file \texttt{SubNLogN/\allowbreak Challenge.lean}. Its core definitions are quoted in Section~\ref{sec:lean}.\looseness=-1

\paragraph{Remarks.}
Two remarks clarify how Theorems~\ref{thm:main} and~\ref{thm:lean} relate to prior work and to the exponent.
\begin{itemize}[leftmargin=*,itemsep=1.5pt,topsep=2pt]
  \item \emph{What is removed and what is borrowed.}
    Our proof uses no result from~\cite{HvdH21} and discards its prime-size grids and polynomial-ring layers. The written proof of Theorem~\ref{thm:main} borrows from~\cite{OAI26} only two self-contained black boxes: its linear-time tape-stream operations and its two fixed linear networks at ground-set parameter $h=100$ with their finite algebraic certificates, stated in Section~\ref{sec:prelim-oai}. The Lean~4 proof of Theorem~\ref{thm:lean} assumes neither black box; it constructs and verifies the tape operations, both $h=100$ networks, and their certificates from scratch.

  \item \emph{The exponent is a by-product.}
    Simplifying onto one power-of-two grid lets us present one direct route, with the larger $\kappa$ emerging as a by-product. Because the half-sum and half-difference Fourier step also permutes address-bit positions bit-exactly, removing the layout bottleneck raises $\kappa$ to $2^{-56}$ under the constants in~\cite{OAI26}. Retuning with the certified constants of Appendix~\ref{app:parameter-transfer} reaches $\kappa=2^{-46}$; Table~\ref{tab:kappa} in Section~\ref{sec:overview-fft} attributes the gain over $2^{-182}$. These exponents are purely asymptotic, as the sub-$n\log n$ regime begins only at astronomical sizes; see Sections~\ref{sec:overview} and~\ref{sec:discussion}.\looseness=-1
\end{itemize}

\paragraph{Use of AI tools.}
This work was developed with substantial assistance from AI agents (primarily Gemini~4), working under the direction of the authors. That assistance spanned the algorithm, the proofs in this paper, the exact-arithmetic reference implementation, and the Lean~4 formalization. The authors have checked the written proofs in this paper. No author has checked the formal Lean proof line by line. Rather, the correctness of the formal development rests on kernel verification by Lean~4 (replayed through the Lean kernel with \texttt{leanchecker} and independently checked with \texttt{comparator} and \texttt{nanoda}), and the authors have checked the formal statement in \texttt{SubNLogN/\allowbreak Challenge.lean} and its correspondence to Theorem~\ref{thm:main}. Agreement among AI systems was used only as a working heuristic during exploration, never as mathematical evidence: every claim in this paper is established by a complete written proof, an explicit citation to the literature, or a kernel-checked formal theorem.

\subsection{Organization}\label{sec:intro-roadmap}

Section~\ref{sec:overview} explains how the algorithm works; the full statements and proofs begin in Section~\ref{sec:prelim}. Section~\ref{sec:prelim} fixes notation, the Turing machine and fixed-point conventions, discrete Fourier transform identities, and the two black-box network facts. Section~\ref{sec:bridge} builds the multidimensional bridge that smooths each spread axis, undoes that smoothing to working precision, and fixes the precision budget. A reader who accepts that precision contract may skip its error proofs on a first reading and proceed to Section~\ref{sec:walsh}, which uses pairwise half-sum and half-difference steps to permute address bits bit-exactly.\looseness=-1

Section~\ref{sec:direct-fft} uses those address permutations to compute the multidimensional Fourier transform directly over the complex numbers. Section~\ref{sec:main} combines the bridge and transform costs, fixes the schedule exponents, and proves Theorem~\ref{thm:main}. Section~\ref{sec:lean} describes the Lean~4 formalization of Theorem~\ref{thm:lean}; its formal statement, trust boundary, and kernel-replay audit can be read independently of Sections~\ref{sec:bridge}--\ref{sec:main}. Section~\ref{sec:discussion} examines where the exponent comes from, why activation sizes are astronomical, what became simpler, and what remains open.

The appendices hold supporting material that may be skipped on a first reading. Appendix~\ref{app:parameter-transfer} re-derives the address-interchange and parallel Fourier-step bounds with general parameters and certifies the $h=100$ network exponents. Appendix~\ref{app:deferred} proves two impossibility bounds for alternative smoothing kernels together with series bounds for precomputed tables. Appendix~\ref{app:bluestein} analyzes the alternative polynomial-ring route and extra parameter rows for comparison. Appendix~\ref{app:oai} records the notation map to prior work, and Appendix~\ref{app:impl} summarizes the exact-arithmetic Python implementation.\looseness=-1
